\section{Gauge and the Standard Model}

Combining the coin with the tone gives a grand-unified gauge structure that breaks to the Standard Model, and it
comes with two sets of quantitative predictions. The gauge group is not chosen. It is forced by the substrate's
own content once the tone is admitted as a fifth axis, and the symmetry breaking that yields the Standard Model
is automatic rather than tuned.

A few terms first. A \emph{gauge} field is the kind of field that carries a force, the electromagnetic field
being the familiar one. The word gauge names a redundancy, since there are many equally valid ways to write down
the same physical field, and that freedom of description is the force's symmetry. A grand unified theory, such as
$\mathrm{SO}(10)$ here, is a single larger symmetry that holds the Standard Model's three forces as pieces,
joined into one at high energy.

\subsection{$\mathrm{SO}(10)$ from coin and tone}

The coin's $D_4$ root system, $\mathrm{SO}(8)$, together with the ternary tone treated as an extra weight, builds
the $D_5 = \mathrm{SO}(10)$ root system. The $16$-dimensional spinor of $\mathrm{SO}(10)$ is exactly one
generation, with $16 = 8s + 8c$ split by the tone sign, the $8s$ carrying tone $+$ and the $8c$ tone $-$. So the
gauge group and one full family of matter both come from coin and tone.

\begin{theorem}[$\mathrm{SO}(10)$ from coin and tone]
The $D_4$ root system of the coin, extended by the ternary tone as a fifth axis, is the $D_5 = \mathrm{SO}(10)$
root system. The $\mathrm{SO}(10)$ spinor $16 = 8s + 8c$, split by tone sign, is exactly one generation of
matter. The grand-unified gauge group and one family arise together from the substrate's coin and tone.
\end{theorem}

The tone is needed, not a free add-on. The coin alone, $\mathrm{SO}(8) = D_4$, provably cannot hold the Standard
Model, and the reason is an orthogonality obstruction. The group $\mathrm{SU}(3) \times \mathrm{SU}(2)$ embeds in
$\mathrm{SO}(N)$ only if there is an $A_2$ root system with a root orthogonal to all of it. Every $D_4$ root has
exactly two non-zero entries, so no $D_4$ root is orthogonal to a full $A_2$. The Standard Model does not fit in
$\mathrm{SO}(8)$.

\begin{proposition}[The orthogonality obstruction]
The Standard Model gauge group does not embed in $\mathrm{SO}(8) = D_4$. An embedding of
$\mathrm{SU}(3) \times \mathrm{SU}(2)$ requires an $A_2$ root system with a root orthogonal to it, but every
$D_4$ root has exactly two non-zero entries and so cannot be orthogonal to a full $A_2$. Adding the ternary tone
as a fifth axis removes the obstruction, since $D_4$'s $24$ roots plus the tone sign give the $40$ roots of
$D_5 = \mathrm{SO}(10)$, the $16$ new roots being $8v$ tensored with the tone sign.
\end{proposition}

The tone is forced by the obstruction, not chosen. By root counting, $D_4$'s $24$ roots plus the tone sign give
$D_5 = \mathrm{SO}(10)$'s $40$ roots, and the $16$ extra roots are precisely $8v$ tensored with the tone sign, so
the tone promotes $\mathrm{SO}(8)$ to $\mathrm{SO}(10)$. A further bonus comes for free. The $16$ of
$\mathrm{SO}(10)$ is anomaly-free, so anomaly cancellation and charge quantization are automatic consequences of
the spinor structure rather than separate conditions to be imposed.

\subsection{Breaking to the Standard Model}

A self-condensate, playing the role of the grand-unified Higgs, breaks
$\mathrm{SO}(10) \to \mathrm{SU}(5) \to \mathrm{SU}(3) \times \mathrm{SU}(2) \times \mathrm{U}(1)$ along the
Georgi-Glashow chain \cite{georgiglashow1974, fritzschminkowski1975}, with $16 = 1 + 10 + \overline{5}$ under
$\mathrm{SU}(5)$, the complete matter content of one generation, quarks, leptons, and a right-handed neutrino.
Any self-condensate gives $\mathrm{SU}(5)$ with $20$ unbroken roots, so the breaking is automatic rather than
fine-tuned. Triality gives three generations.

\begin{proposition}[Automatic breaking, and the self as Higgs]
A self-condensate breaks $\mathrm{SO}(10)$ to $\mathrm{SU}(5)$ for any choice of condensate, leaving $20$
unbroken roots, and the chain continues to
$\mathrm{SU}(3) \times \mathrm{SU}(2) \times \mathrm{U}(1)$. Under $\mathrm{SU}(5)$ the generation decomposes as
$16 = 1 + 10 + \overline{5}$. The breaking field is itself a self-condensate, so the matter sector and the
symmetry-breaking sector are the same object.
\end{proposition}

The clean conceptual point is that the breaking field is a self-condensate, so the self is the Higgs. The matter
sector and the symmetry-breaking sector are not two things but one.

\subsection{The weak mixing angle}

Summing over the fermions of one generation gives the weak mixing angle at unification,
\[
\sin^2\theta_W = \frac{\sum T_3^2}{\sum Q^2} = \frac{3}{8}
\]
exactly at the unification scale \cite{georgiquinnweinberg1974}. Running it down with standard renormalization
group flow, it reaches the measured $0.231$ for a supersymmetric spectrum, where the three gauge couplings unify
cleanly near $10^{16}$ GeV, or about $0.21$ for the bare Standard Model.

\begin{theorem}[$\sin^2\theta_W = 3/8$ at unification]
Summed over one generation, $\sin^2\theta_W = \sum T_3^2 / \sum Q^2 = 3/8$ exactly at the unification scale. Run
to the weak scale it gives the observed $0.231$ for a supersymmetric spectrum and about $0.21$ for the bare
Standard Model.
\end{theorem}

So the prediction is confirmed against experiment, and it additionally points to low-energy supersymmetry, since
the cleaner agreement comes with the supersymmetric running and the matching gauge-coupling unification.

\subsection{The grand-unified mass relations}

Three mass relations follow at the unification scale. The first is $b$-$\tau$ unification, $m_b = m_\tau$ at the
grand-unified scale from the shared Yukawa coupling, running to $m_b / m_\tau \approx 2.3$ at $M_Z$, against the
observed $2.35$. The second is the determinant relation
$\det(M_{\text{charged lepton}}) = \det(M_{\text{down}})$, which follows from the discrete fact that the trace of
the Yukawa over the $16$ vanishes. The third is the Georgi-Jarlskog factor-three relations
\cite{georgijarlskog1979}, the weakest of the set and the most model-dependent.

\subsection{The non-abelian gauge field}

The gauge field is genuine, and this was verified directly. Using $\mathrm{SO}(3)$ link matrices as the
prototype, identical in structure to the $\mathrm{SO}(10)$ case, the Wilson loop, the trace of the
ordered product of links around a loop, is gauge-invariant \cite{wilson1974}, the holonomy is non-trivial, which
is curvature, and the link matrices do not commute, which is the non-abelian content.

\begin{result}[A non-abelian gauge field on the substrate]
Built from $\mathrm{SO}(3)$ link matrices on the substrate, the Wilson loop is gauge-invariant, its holonomy is
non-trivial, and the links do not commute \cite{pollard2026vibetest}. The gauge field is genuine, curved, and
non-abelian, with the $\mathrm{SO}(10)$ case identical in structure.
\end{result}

The symmetric collision that makes $\mathrm{SO}(10)$ the local symmetry is forced in the infrared by triality.
The cleaner principle underneath is that a local conservation law is a Gauss law, so the rule's exact charge
conservation gives $\mathrm{U}(1)$ electromagnetism for free, taken up where electromagnetism is treated, and a
non-abelian gauge group reduces to the collision's symmetry. Gauging $\mathrm{SO}(10)$ is then just the existence
of an $\mathrm{SO}(10)$-symmetric collision, which the substrate supplies.


The whole gauge and matter sector, together with the predictions $\sin^2\theta_W = 3/8$ and the mass relations,
comes from coin and tone, with the Standard Model emerging by automatic breaking and with anomaly cancellation
included for free.
